\chapter{Node Operations}

\section{The Node as Infrastructure}

A Connector node is not just a process — it is a governed runtime that manages
agents, memory, and protocol sessions with the same rigour as a database server
or a message broker. This chapter covers what happens from the moment
\texttt{connector-node} starts to the moment it shuts down gracefully, and
every operational concern in between.

\section{The 12-Stage Boot Sequence}

Connector nodes boot through a deterministic 12-stage pipeline. Each stage
either succeeds and advances to the next, or fails and halts the boot with a
clear diagnostic. There is no partial initialisation and no silent degraded
mode — the node is either fully ready or it does not serve traffic.

\begin{figure}[H]
\centering
\begin{tikzpicture}[
  st/.style={rectangle, rounded corners=3pt, draw=#1!70, fill=#1!8,
             text width=2.85cm, align=center, minimum height=0.85cm,
             font=\scriptsize\bfseries, inner sep=4pt},
  node distance=0.32cm
]
%% Row 1 (stages 1-6)
\node[st=CInfra]   (s1)  at (0,2.0)     {1\\IDENTITY};
\node[st=CInfra]   (s2)  at (3.05,2.0)  {2\\CONFIG};
\node[st=CInfra]   (s3)  at (6.1,2.0)   {3\\SECRETS};
\node[st=CKernel]  (s4)  at (9.15,2.0)  {4\\STORAGE};
\node[st=CKernel]  (s5)  at (12.2,2.0)  {5\\KERNEL};
\node[st=CLogic]   (s6)  at (12.2,0.8)  {6\\POLICIES};

%% Row 2 (stages 7-12, right to left so arrow flows naturally)
\node[st=CLogic]   (s7)  at (9.15,0.8)  {7\\SCHEDULER};
\node[st=CExec]    (s8)  at (6.1,0.8)   {8\\CAPABILITIES};
\node[st=CExec]    (s9)  at (3.05,0.8)  {9\\RESTORE};
\node[st=CProto]   (s10) at (0,0.8)     {10\\SERVICES};
\node[st=CDX]      (s11) at (0,-0.45)   {11\\ACCESS};
\node[st=CBiz]     (s12) at (3.05,-0.45){12\\READY};

%% Row 1 arrows
\foreach \a/\b in {s1/s2,s2/s3,s3/s4,s4/s5}
  \draw[arr=gray!60] (\a) -- (\b);
\draw[arr=gray!60] (s5.south) -- (s6.north);

%% Row 2 arrows
\foreach \a/\b in {s6/s7,s7/s8,s8/s9,s9/s10}
  \draw[arr=gray!60] (\a) -- (\b);
\draw[arr=gray!60] (s10.south) -- (s11.north);
\draw[arr=gray!60] (s11.east)  -- (s12.west);

%% Label
\node[above=0.12cm of s1, font=\scriptsize\itshape, text=gray]{boot start};
\node[above=0.12cm of s12, font=\scriptsize\itshape, text=CBiz]{sd\_notify READY};

\end{tikzpicture}
\caption{Connector node 12-stage boot pipeline — any stage failure halts the boot}
\label{fig:boot-pipeline}
\end{figure}

\begin{table}[H]
\centering
\renewcommand{\arraystretch}{1.35}
\begin{tabularx}{\textwidth}{@{}clX@{}}
\toprule
\textbf{\#} & \textbf{Stage} & \textbf{What happens} \\
\midrule
1  & IDENTITY     & Load or generate node DID; load Ed25519 signing keys \\
2  & CONFIG       & Parse \texttt{config.toml}, env vars, CLI flags; resolve hierarchy \\
3  & SECRETS      & Load API tokens, LLM keys, database credentials from vault/file \\
4  & STORAGE      & Connect to storage backend (PostgreSQL/SQLite); verify schema \\
5  & KERNEL       & Boot Memory Kernel; verify CID chain integrity; load namespace index \\
6  & POLICIES     & Load governance predicates and the 11-layer predicate chain \\
7  & SCHEDULER    & Initialise agent scheduler; allocate resource pools \\
8  & CAPABILITIES & Register tool providers: built-in, MCP servers, CNP ports \\
9  & RESTORE      & Reload agents and sessions from last checkpoint; resume execution \\
10 & SERVICES     & Start HTTP API, WebSocket, CNP router, A2A and MCP bridges \\
11 & ACCESS       & Enable local dashboard; enable API key authentication \\
12 & READY        & Send \texttt{sd\_notify READY}; start health probe endpoints \\
\bottomrule
\end{tabularx}
\caption{12-stage boot — description of each stage}
\end{table}

Each stage writes a boot-progress log line with its name, status, and elapsed
time. The terminal output follows a fixed format:
\texttt{[N/12] STAGE-NAME  OK (0.42s)}, which makes boot logs scannable by eye
and parseable by automation.

\section{systemd Integration}

Connector nodes are designed as systemd services. The key integration points
are:

\begin{itemize}[leftmargin=1.5em, itemsep=2pt]
  \item \textbf{\texttt{Type=notify}} — systemd waits for \texttt{READY=1}
        before marking the service as started, preventing premature traffic
        routing.
  \item \textbf{\texttt{WatchdogSec=60s}} — the node must send a watchdog
        heartbeat every 60 seconds or systemd kills and restarts it.
  \item \textbf{\texttt{Restart=on-failure}} — automatic restart on any
        non-zero exit; agents resume from checkpoint on next boot.
\end{itemize}

\section{Health Probes}

Connector exposes three health probe endpoints compatible with both Kubernetes
and traditional load balancers:

\begin{table}[H]
\centering
\renewcommand{\arraystretch}{1.45}
\begin{tabularx}{\textwidth}{@{}llXl@{}}
\toprule
\textbf{Probe} & \textbf{Endpoint} & \textbf{What it checks} & \textbf{Failure means} \\
\midrule
Liveness  & \texttt{GET /health/live}    &
  Node event loop is not deadlocked &
  Kill and restart immediately \\
Readiness & \texttt{GET /health/ready}   &
  Kernel, scheduler, and storage are all operational &
  Remove from load balancer; stop sending traffic \\
Startup   & \texttt{GET /health/startup} &
  All 12 boot stages have completed &
  Do not start other probes yet; give more boot time \\
\bottomrule
\end{tabularx}
\caption{Health probe endpoints — Kubernetes and load-balancer compatible}
\end{table}

\begin{insightbox}{Probe ordering in Kubernetes}
Configure the startup probe first with a generous \texttt{failureThreshold}
(e.g., 30 attempts × 10s = 5 minutes). Only after startup succeeds do the
liveness and readiness probes activate. This prevents Kubernetes from killing
a slow-starting node during a large checkpoint restore.
\end{insightbox}

\section{Agent Deployment}

Agents are deployed from a TOML manifest file. The manifest declares the
agent's identity, runtime requirements, tool capabilities, governance
predicates, and resource budget.

\subsection{Deployment Flow}

\begin{figure}[H]
\centering
\begin{tikzpicture}[
  st/.style={rectangle, rounded corners=3pt, draw=CExec!60, fill=CExec!8,
             text width=2.5cm, align=center, minimum height=0.85cm,
             font=\footnotesize\bfseries},
  node distance=0.7cm
]
\node[st] (parse)   {Parse\\manifest};
\node[st, right=of parse]  (verify) {Verify\\capabilities};
\node[st, right=of verify] (gov)    {Compile\\governance};
\node[st, right=of gov]    (alloc)  {Allocate\\resources};
\node[st, right=of alloc]  (reg)    {Register\\in kernel};
\node[st, right=of reg]    (start)  {Start\\agent};
\node[st, right=of start]  (hc)    {Health\\check};

\foreach \a/\b in {parse/verify,verify/gov,gov/alloc,alloc/reg,reg/start,start/hc}
  \draw[arr=CExec!60] (\a) -- (\b);

\node[below=0.25cm of gov, font=\scriptsize\itshape, text=gray]
  {predicates → executable form};
\node[below=0.25cm of start, font=\scriptsize\itshape, text=gray]
  {DID assigned};

\end{tikzpicture}
\caption{Agent deployment pipeline — manifest to running agent}
\end{figure}

\subsection{Agent Manifest Structure}

The manifest is a typed, validated TOML document. Four top-level sections
define the agent:

\begin{table}[H]
\centering
\renewcommand{\arraystretch}{1.4}
\begin{tabularx}{\textwidth}{@{}lX@{}}
\toprule
\textbf{Section} & \textbf{What it declares} \\
\midrule
\texttt{[agent]}       & Name, description, version, DID prefix \\
\texttt{[runtime]}     & LLM model, temperature, max tokens, cognitive substrate config \\
\texttt{[capabilities]}& Tool names, memory namespaces, protocol ports to open \\
\texttt{[governance]}  & Predicate names, resource budget (tokens, memory, time, calls) \\
\bottomrule
\end{tabularx}
\caption{Agent manifest sections}
\end{table}

\section{Monitoring and Observability}

\subsection{Prometheus Metrics}

The node exposes a \texttt{/metrics} endpoint in Prometheus text format.
Key metrics exported:

\begin{table}[H]
\centering
\renewcommand{\arraystretch}{1.4}
\rowcolors{2}{gray!5}{white}
\begin{tabular}{@{}ll@{}}
\toprule
\textbf{Metric name} & \textbf{Type and description} \\
\midrule
\texttt{connector\_agents\_running}     & Gauge — active agent count \\
\texttt{connector\_sessions\_total}     & Counter — sessions opened since start \\
\texttt{connector\_memory\_packets}     & Counter — MemPackets written \\
\texttt{connector\_tokens\_consumed}    & Counter — total LLM tokens used \\
\texttt{connector\_predicate\_checks}   & Counter — predicate enforcement calls \\
\texttt{connector\_predicate\_denials}  & Counter — predicate enforcement denials \\
\texttt{connector\_tool\_calls\_total}  & Counter — external tool invocations \\
\texttt{connector\_boot\_duration\_ms}  & Gauge — last boot time in milliseconds \\
\texttt{connector\_receipts\_issued}    & Counter — signed execution receipts \\
\bottomrule
\end{tabular}
\caption{Core Prometheus metrics exposed by the node}
\end{table}

\subsection{Structured Logging}

All node logs are emitted in structured JSON via the \texttt{tracing} framework.
Every log line carries a consistent set of fields: \texttt{timestamp},
\texttt{level}, \texttt{target} (module), \texttt{agent\_did},
\texttt{session\_id}, and \texttt{message}. This makes logs directly
ingestible by Loki, Elasticsearch, Splunk, or any JSON-aware log pipeline.

\subsection{OpenTelemetry Tracing}

Long-running operations (agent execution, contract compilation, checkpoint
restore) emit OpenTelemetry spans. Each span carries the agent DID, session
CID, and predicate enforcement result. Traces can be exported to Jaeger,
Tempo, or any OTLP-compatible collector.

\section{Graceful Shutdown}

Connector nodes shut down in a deterministic eight-phase sequence on receipt
of \texttt{SIGTERM} or \texttt{SIGINT}. No agent is killed mid-execution;
all in-progress sessions are checkpointed before the process exits.

\begin{figure}[H]
\centering
\begin{tikzpicture}[
  ph/.style={rectangle, rounded corners=3pt, draw=CInfra!70, fill=CInfra!8,
             text width=3.2cm, align=center, minimum height=0.9cm,
             font=\scriptsize\bfseries},
  node distance=0.5cm and 0.5cm
]
%% Row 1
\node[ph] (sig)    {Receive\\SIGTERM};
\node[ph, right=of sig]  (drain) {Stop\\accepting};
\node[ph, right=of drain] (inflight) {Drain\\in-flight (30s)};
\node[ph, right=of inflight] (pause) {Checkpoint\\agents};

%% Row 2
\node[ph, below=0.7cm of pause] (sess)  {Close\\sessions};
\node[ph, left=of sess]  (flush) {Flush\\memory};
\node[ph, left=of flush] (svc)   {Shutdown\\services};
\node[ph, left=of svc]   (exit)  {Exit 0};

%% Row 1 arrows
\foreach \a/\b in {sig/drain,drain/inflight,inflight/pause}
  \draw[arr=CInfra!60] (\a) -- (\b);
\draw[arr=CInfra!60] (pause.south) -- (sess.north);

%% Row 2 arrows
\foreach \a/\b in {sess/flush,flush/svc,svc/exit}
  \draw[arr=CInfra!60] (\a) -- (\b);

\node[below=0.18cm of inflight, font=\scriptsize\itshape, text=gray]
  {503 from health probes};
\node[below=0.18cm of flush, font=\scriptsize\itshape, text=gray]
  {all MemPackets persisted};

\end{tikzpicture}
\caption{Graceful shutdown sequence — 8 phases, no data loss}
\label{fig:shutdown}
\end{figure}

\section{Backup and Restore}

The kernel's content-addressed store makes backup straightforward:
a backup is a consistent snapshot of the MemPacket store. Because every
packet has a CID, incremental backups are simple set differences.

\begin{lstlisting}[language=bash]
connectorctl backup memory --output /backups/$(date +%Y%m%d).tar.gz
connectorctl backup memory --incremental --since 1h --output /backups/inc.tar.gz
connectorctl restore memory --input /backups/20260322.tar.gz
\end{lstlisting}

Agent-level backups capture the agent's full MemPacket namespace and the last
signed checkpoint, making it possible to migrate an agent to a new node with
full history intact.

\section{Disaster Recovery}

\begin{table}[H]
\centering
\renewcommand{\arraystretch}{1.5}
\rowcolors{2}{gray!5}{white}
\begin{tabularx}{\textwidth}{@{}lXl@{}}
\toprule
\textbf{Scenario} & \textbf{Recovery procedure} & \textbf{Data loss} \\
\midrule
Node crash            & Restart node; agents resume from last checkpoint on stage 9 & None \\
Disk corruption       & Restore from backup; replay audit receipt chain & Up to last backup \\
Database failure      & Switch to replica (if configured); restore from WAL & Seconds \\
Network partition     & Agents pause; session receipts held in local queue; resume on reconnect & None \\
License expiry        & Agents pause gracefully; resume after license renewal via \texttt{connectorctl activate} & None \\
Kernel integrity fail & Stage 5 halts boot; operator restores from backup or runs \texttt{connectorctl repair} & None \\
\bottomrule
\end{tabularx}
\caption{Disaster recovery scenarios and procedures}
\end{table}

\begin{table}[H]
\centering
\renewcommand{\arraystretch}{1.4}
\begin{tabularx}{\textwidth}{@{}lXX@{}}
\toprule
\textbf{Metric} & \textbf{Single-node} & \textbf{Cluster (3+ nodes)} \\
\midrule
RTO (Recovery Time Objective) & $\leq 5$ minutes      & $\leq 30$ seconds \\
RPO (Recovery Point Objective)& Last checkpoint ($<$1 min) & Near-zero (replica sync) \\
\bottomrule
\end{tabularx}
\caption{RTO and RPO targets by deployment topology}
\end{table}

\section{Deployment Topologies}

\begin{figure}[H]
\centering
\begin{tikzpicture}[
  node_/.style={rectangle, rounded corners=4pt, draw=CInfra!70, fill=CInfra!8,
                text width=2.4cm, align=center, minimum height=1.1cm,
                font=\small\bfseries},
  cluster/.style={rectangle, rounded corners=6pt, draw=CKernel!40, fill=CKernel!4,
                  dashed, inner sep=8pt},
  node distance=1.2cm
]
%% Single-node
\node[node_] (sn) at (-5, 0) {Single\\Node};
\node[below=0.15cm of sn, font=\footnotesize\itshape, text=gray]
  {Dev / small prod};

%% Cluster
\node[node_] (n1) at (0,  0.9) {Node 1\\(Primary)};
\node[node_] (n2) at (2.0, 0)  {Node 2\\(Replica)};
\node[node_] (n3) at (0, -0.9) {Node 3\\(Replica)};

\begin{scope}[on background layer]
  \node[cluster, fit=(n1)(n2)(n3)] (cl) {};
\end{scope}

\node[above=0.1cm of cl.north, font=\footnotesize\itshape, text=gray]
  {Cluster — shared kernel store};

\draw[darr=CKernel!50] (n1) -- (n2);
\draw[darr=CKernel!50] (n1) -- (n3);
\draw[darr=CKernel!50, dashed] (n2) -- (n3);

%% Cloud / hybrid
\node[node_] (hybrid) at (6, 0) {Hosted\\SaaS Node};
\node[below=0.15cm of hybrid, font=\footnotesize\itshape, text=gray]
  {Managed by Connector};

\node[above=0.3cm of sn.north, font=\small\bfseries, text=CInfra]
  {Local};
\node[above=0.3cm of hybrid.north, font=\small\bfseries, text=CInfra]
  {Cloud};

\end{tikzpicture}
\caption{Supported deployment topologies — single-node, cluster, and hosted SaaS}
\label{fig:topologies}
\end{figure}

\begin{table}[H]
\centering
\renewcommand{\arraystretch}{1.4}
\begin{tabularx}{\textwidth}{@{}lXXX@{}}
\toprule
 & \textbf{Single-node} & \textbf{Cluster} & \textbf{Hosted SaaS} \\
\midrule
\textbf{License tier}    & Free–Enterprise   & Enterprise       & Any \\
\textbf{HA}              & No                & Yes              & Yes \\
\textbf{Data location}   & Customer premises & Customer premises & Connector-managed \\
\textbf{Management}      & Self-managed      & Self-managed      & Fully managed \\
\textbf{Best for}        & Dev, small prod   & Large prod, HA    & No ops team \\
\bottomrule
\end{tabularx}
\caption{Deployment topology comparison}
\end{table}

